% Imporditud: tools/import_eksperimendid.py
% Allikas: Eesti füüsikaolümpiaadi eksperimendi ülesannete
%   kogu (Rael Kalda, 2023), ülesanne Ü13, lahendus L13.
\setAuthor{EFO žürii}
\setRound{piirkonnavoor}
\setYear{2006}
\setNumber{P E1}
\setDifficulty{1}
\setTopic{vedelike mehaanika}
\setVahendid{anum veega, 4 kustutuskummi, nööpnõelad, dünamomeeter}
\setPunktid{8}
\setAllikas{Eksperimendi kogu Ü13, lahendus lk 38}
\setLahendusseis{toores}

\prob{Kustutuskumm}
Määrake kustutuskummi tihedus. Vee tihedus $\rho$ = 1,0 g/cm$^3$.

\hint

\solu
Katse idee --- 3 p. Katsekeha kaal õhus on F = mg ja vees P = F $- \rho$vgV (1 p). Arvestades, et V = m/$\rho$, saame

( ) m $\Rightarrow$ P = F 1 $- \rho$v $\cdot$ (1,5 p) P = F $- \rho$vgV $\Rightarrow$ P = F $- \rho$vg $\rho$ $\rho$

Avaldame viimasest võrrandist katsekeha tiheduse $\rho$, saame

$\rho$ = F $\rho$v . (0,5 p) F $-$P

Mõõtmised --- 1 p, lõpptulemuse arvutamine --- 1 p.
\probend
